La implementación computacional validó con éxito el modelo XGBoost en la sección transversal de prueba de IBM Telco ($N=1,409$). Desde el punto de vista algorítmico estándar, el clasificador capturó correctamente las señales multivariadas de desgaste reportando un rendimiento elevado, lo que fundamenta sólidas capacidades deductivas como motor estadístico principal.

Sin embargo, el eje central de nuestros resultados emerge al acoplar estas predicciones estáticas en el motor dinámico de Inteligencia de Negocios (BI). Durante la simulación de intervención interactiva ($C_R = 50$ USD como costo de retención e $I_S = 300$ USD de rescate), identificamos que optimizar la curva de decisión empujando un umbral ($\tau$) conservador sobre los perfiles propensos previene el desgaste comercial de campañas masivas ciegas. El cruce entre el número de Falsos Positivos no intervenidos (eludiendo la purga por $C_R$) y de Verdaderos Positivos reconvertidos, dictaminó una rentabilidad neta teórica escalable bajo la matriz de confusión interactiva.

Adicionalmente, el volcado gráfico de los valores de SHAP locales proveyó a la interfaz BI de una anatomía del riesgo unívoca. Esto facultó operativamente a los tomadores de decisión para diagnosticar a los clientes evaluados de forma unitaria, separando inmediatamente entre aquellos individuos empujados al abandono por factores financieros directos (\textit{MonthlyCharges}) en contraposición a aquellos sin lazos de lealtad estructurales técnicos (ausencia de \textit{TechSupport} o \textit{OnlineSecurity}).
